\documentclass[11pt]{article}
\usepackage[a4paper,margin=1in]{geometry}
\usepackage[T1]{fontenc}
\usepackage[utf8]{inputenc}
\usepackage{lmodern,microtype}
\usepackage{amsmath,amssymb,amsthm,mathtools}
\usepackage{enumitem,xcolor,hyperref}
\usepackage[nameinlink,capitalise,noabbrev]{cleveref}
\hypersetup{colorlinks=true,linkcolor=blue!60!black,citecolor=blue!60!black,urlcolor=blue!60!black}

\newtheorem{definition}{Definition}[section]
\newtheorem{assumption}[definition]{Assumption}
\newtheorem{theorem}[definition]{Theorem}
\newtheorem{proposition}[definition]{Proposition}
\newtheorem{corollary}[definition]{Corollary}
\newtheorem{conjecture}[definition]{Conjecture}
\newtheorem{remark}[definition]{Remark}

\newcommand{\C}{\mathbb C}
\newcommand{\Oterm}{\mathcal O}
\newcommand{\D}{\mathrm D}
\newcommand{\Q}{\mathcal Q}
\newcommand{\norm}[1]{\left\lVert #1\right\rVert}
\DeclareMathOperator{\Sym}{Sym}
\DeclareMathOperator{\spanop}{span}

\title{Finite-Probe Inverse-Jet Reconstruction for Pandrosion\\[3pt]
\large Symmetric probe stencils, tensor recovery, and derivative-free chart flattening}
\author{Ivan Besevic\\Independent Mathematical Research}
\date{May 2026}

\begin{document}
\maketitle

\begin{abstract}
The all-order tensor inverse-jet theorem identifies the optimal local chart for a regular zero of an analytic system, but it is stated in terms of Taylor tensors. A derivative-free Pandrosion engine should recover those tensors from finite probes. This paper formulates and proves the local consistency statement behind that goal. For an analytic map \(G:\C^n\to\C^n\) near a regular zero \(\rho\), a fixed total-degree unisolvent probe cloud, evaluated at radius \(h\), recovers the homogeneous Taylor coefficients \(G_q\) with error \(\Oterm(h)\). The reconstruction is done by one scaled polynomial fit, not by recursively subtracting lower-order \(\Oterm(h)\) approximations; this avoids the loss of powers of \(h\) that would otherwise contaminate higher tensors. Applying the inverse-jet recursion to the recovered tensors gives estimated chart coefficients \(\widehat B_q(h)\) satisfying
\[
        \widehat B_q(h)\longrightarrow B_q,\qquad 1\le q\le m,
\]
where \(B_q\) are the exact all-order inverse-jet coefficients. Consequently the reconstructed chart \(\widehat\phi_m\) is asymptotically \(m\)-flat:
\[
        G(\widehat\phi_m(u))=u+\Oterm_h(\norm{u}^{m+1})+o(1)
\]
on bounded chart coordinates. The paper closes with a finite-slope engine conjecture: stable Pandrosion probes with rank filtering and equal-value avoidance implement this reconstruction asymptotically without explicit derivative access.
\end{abstract}

\paragraph{Scope and status.}
The main results are local consistency theorems for finite probe reconstruction. They do not claim global convergence of a solver; they supply the derivative-free bridge between tensor inverse jets and practical Pandrosion chart engines.

\section{Problem}

The all-order inverse-jet chart at a regular zero is
\[
        \phi_m(u)=\rho+\sum_{q=1}^m B_q[u,\ldots,u],
\]
chosen so that
\[
        G(\phi_m(u))=u+\Oterm(\norm{u}^{m+1}).
\]
The coefficients \(B_q\) are computable from the Taylor tensors of \(G\). A Pandrosion engine, however, observes values of \(G\) and finite slopes, not exact derivatives. We therefore ask:
\begin{quote}
Can the inverse-jet coefficients be reconstructed from finite probe evaluations as the probe radius tends to zero?
\end{quote}


\section{Probe stencils}

Let \(H_q\) denote the finite-dimensional vector space of homogeneous polynomial maps \(P:\C^n\to\C^n\) of degree \(q\), and let
\[
        \Pi_m^0=H_1\oplus\cdots\oplus H_m
\]
be the space of \(\C^n\)-valued polynomials of total degree at most \(m\) with zero constant term.

\begin{definition}[Homogeneous unisolvency]
A finite set of directions \(V=\{v_1,\ldots,v_N\}\subset\C^n\) is homogeneous-unisolvent through degree \(m\) if, for every \(1\le q\le m\), the evaluation map
\[
        H_q\to(\C^n)^N,\qquad P\mapsto(P(v_1),\ldots,P(v_N))
\]
is injective.
\end{definition}

\begin{definition}[Total-degree probe cloud]
A finite probe cloud \(Y=\{y_1,\ldots,y_M\}\subset\C^n\) is total-degree unisolvent through degree \(m\) if the evaluation map
\[
        \Pi_m^0\to(\C^n)^M,
        \qquad P\mapsto(P(y_1),\ldots,P(y_M))
\]
is injective.
Equivalently, a polynomial in \(\Pi_m^0\) that vanishes on all points of \(Y\) must be zero.
\end{definition}

\begin{remark}
Generic probe clouds are total-degree unisolvent once the number of nodes is at least \(\dim \Pi_m^0\). In implementation, one may use an overdetermined cloud and a fixed least-squares reconstruction operator. Homogeneous-unisolvent directions are still useful, but isolating the degree-\(q\) tensor from a single radius requires a total-degree fit or an equivalent multi-radius finite-difference formula.
\end{remark}

Write the Taylor expansion at a regular zero \(\rho\) as
\[
        G(\rho+w)=\sum_{q\ge1}G_q[w,\ldots,w],
        \qquad G_1=\D G(\rho).
\]

\section{Homogeneous tensor recovery}

\begin{theorem}[Finite-probe Taylor tensor recovery]\label{thm:tensor-recovery}
Let \(G\) be analytic near \(\rho\), and let \(Y\) be total-degree unisolvent through degree \(m\). Let
\[
        E_Y:\Pi_m^0\to(\C^n)^M,
        \qquad E_Y(P)=(P(y_1),\ldots,P(y_M)),
\]
and choose once and for all a bounded linear reconstruction operator
\(L_Y:(\C^n)^M\to\Pi_m^0\) such that \(L_YE_Y=\mathrm{id}_{\Pi_m^0}\).
For \(h>0\), define
\[
        \widehat P_h(y)=\sum_{q=1}^m h^{q-1}\widehat G_q(h)[y,\ldots,y]
        =L_Y\bigl(h^{-1}G(\rho+h y_1),\ldots,h^{-1}G(\rho+h y_M)\bigr).
\]
When \(E_Y\) is square this is ordinary interpolation; when the cloud is overdetermined it is a fixed least-squares-type reconstruction. Then, for each \(1\le q\le m\),
\[
        \widehat G_q(h)\to G_q
        \quad\text{as }h\to0.
\]
If the evaluations of \(G\) are exact, the error is \(\Oterm(h)\).
\end{theorem}

\begin{proof}
Taylor expansion gives, uniformly for \(y\) in the fixed finite cloud \(Y\),
\[
        h^{-1}G(\rho+h y)
        =
        \sum_{q=1}^{m}h^{q-1}G_q[y,\ldots,y]
        +\Oterm(h^m).
\]
Let
\[
        P_h^*(y)=\sum_{q=1}^{m}h^{q-1}G_q[y,\ldots,y]\in\Pi_m^0.
\]
Thus the scaled data differ from \(E_Y(P_h^*)\) by \(\Oterm(h^m)\). Applying the fixed bounded reconstruction operator \(L_Y\) gives \(\widehat P_h-P_h^*=\Oterm(h^m)\) in coefficient norm. In particular, the degree-\(q\) homogeneous coefficient satisfies
\[
        h^{q-1}\bigl(\widehat G_q(h)-G_q\bigr)=\Oterm(h^m).
\]
Hence \(\widehat G_q(h)-G_q=\Oterm(h^{m-q+1})\), which is at least \(\Oterm(h)\) for \(q\le m\). This proves the claim.
\end{proof}

\begin{remark}[Why recursive subtraction is avoided]
If one first reconstructs \(\widehat G_\ell-G_\ell=\Oterm(h)\) and then subtracts those lower orders before dividing by \(h^q\), the lower-order errors are amplified by factors \(h^{\ell-q}\). For \(\ell<q-1\), this amplification can diverge. The total-degree interpolation in \cref{thm:tensor-recovery} isolates all degrees in one finite-dimensional solve, so no lower-order \(\Oterm(h)\) error is divided by a higher power of \(h\).
\end{remark}

\section{Inverse-jet reconstruction}

The inverse-jet recursion is algebraic. Let \(A=G_1\) and let \(B_q\) be the exact coefficients of the inverse jet. The first coefficient is \(B_1=A^{-1}\). At order \(k\), the coefficient \(B_k\) solves
\[
        A B_k[u,\ldots,u]+R_k(G_2,\ldots,G_k;B_1,\ldots,B_{k-1})(u)=0,
\]
where \(R_k\) is the homogeneous degree-\(k\) residual determined by lower data.

\begin{theorem}[Consistency of finite-probe inverse jets]\label{thm:inverse-consistency}
Assume \(G_1\) is invertible and the hypotheses of \cref{thm:tensor-recovery} hold through degree \(m\). Let \(\widehat B_q(h)\) be obtained by applying the inverse-jet recursion to \(\widehat G_1(h),\ldots,\widehat G_q(h)\). Then
\[
        \widehat B_q(h)\to B_q,\qquad 1\le q\le m.
\]
With exact evaluations, \(\widehat B_q(h)-B_q=\Oterm(h)\) for every fixed \(q\le m\).
\end{theorem}

\begin{proof}
The map \(G_1\mapsto G_1^{-1}\) is continuous near the invertible matrix \(G_1\), so \(\widehat B_1(h)\to B_1\). Suppose the result holds through order \(k-1\). The residual polynomial \(R_k\) is a finite algebraic expression in \(G_2,\ldots,G_k\) and \(B_1,\ldots,B_{k-1}\); therefore its finite-probe estimate converges to the exact residual. Multiplication by \(\widehat G_1(h)^{-1}\) is also continuous. Thus \(\widehat B_k(h)\to B_k\). The \(\Oterm(h)\) rate follows from the same induction because all algebraic maps are locally Lipschitz on bounded finite-dimensional neighborhoods with \(G_1\) bounded away from singularity.
\end{proof}

\begin{corollary}[Asymptotic \(m\)-flatness]\label{cor:mflat}
Let
\[
        \widehat\phi_m^{(h)}(u)=\rho+\sum_{q=1}^m\widehat B_q(h)[u,\ldots,u].
\]
For every fixed bounded set of \(u\)-coordinates,
\[
        G(\widehat\phi_m^{(h)}(u))
        =
        u+\Oterm(\norm{u}^{m+1})+\Oterm(h)\sum_{q=1}^m\norm{u}^q.
\]
In particular, \(\widehat\phi_m^{(h)}\) converges coefficientwise to an \(m\)-flat inverse-jet chart.
\end{corollary}

\begin{proof}
Substitute the coefficient convergence from \cref{thm:inverse-consistency} into the finite Taylor expansion of \(G\). The exact inverse jet cancels homogeneous terms through degree \(m\), and the coefficient error contributes the stated lower-order perturbation.
\end{proof}

\section{Noise and conditioning}

\begin{proposition}[Stable overdetermined reconstruction]\label{prop:noise}
Suppose the total-degree reconstruction operator has condition number at most \(\kappa_Y\), and suppose probe evaluations have additive error at most \(\varepsilon_h\). Then the recovered \(q\)-th Taylor tensor has an additional error bounded by
\[
        \Oterm(\kappa_Y\,\varepsilon_h h^{-q})
\]
at order \(q\). Consequently consistent order-\(m\) reconstruction requires \(\varepsilon_h=o(h^m)\) for the highest tensor.
\end{proposition}

\begin{proof}
The scaled data \(h^{-1}G(\rho+h y_j)\) inherit additive error at most \(\varepsilon_h h^{-1}\). The fixed polynomial reconstruction amplifies this by at most \(\kappa_Y\). The degree-\(q\) coefficient of \(\widehat P_h\) is \(h^{q-1}\widehat G_q(h)\), so division by \(h^{q-1}\) gives an additional tensor error of size \(\Oterm(\kappa_Y\varepsilon_h h^{-q})\). The highest recovered degree imposes the strongest scaling requirement.
\end{proof}

\section{Pandrosion engine conjecture}

\begin{conjecture}[Derivative-free inverse-jet reconstruction by stable probes]
A Pandrosion engine using symmetric total-degree probe clouds, finite-slope rank filtering, equal-value pole avoidance, and shrinking probe radius \(h\) reconstructs the all-order inverse-jet coefficients through degree \(m\) whenever the probe cloud remains uniformly unisolvent and the effective evaluation noise is \(o(h^m)\).
\end{conjecture}

\begin{conjecture}[Reconstructed chart multiplier law]
Under the hypotheses above, the charted anchored Pandrosion map built from the reconstructed chart satisfies
\[
        \norm{D P_{G\circ\widehat\phi_m^{(h)},a}(0)}
        \le
        C\norm{a}^m+o(1)
\]
as \(h\to0\), uniformly for anchors in a fixed small chart ball.
\end{conjecture}

\section{Conclusion}

The all-order inverse jet is not merely a formal derivative object. It can be recovered from finite probes, order by order, provided the stencil is rich enough and the probe radius is controlled. This makes the \(m\)-flat atlas programme algorithmic: exact derivatives define the target, while symmetric finite probes supply a derivative-free path toward it.

\begin{thebibliography}{9}
\bibitem{jets}
I. Besevic, \emph{Higher-Order Jet Cancellation for Anchored Divided-Difference Iterations}, manuscript, May 2026.
\bibitem{tensor}
I. Besevic, \emph{Tensorial Multidimensional Pandrosion}, manuscript, May 2026.
\bibitem{allorder}
I. Besevic, \emph{All-Order Tensor Inverse Jets for Multidimensional Pandrosion}, manuscript, May 2026.
\bibitem{atlas}
I. Besevic, \emph{Uniform \(m\)-Flat Pandrosion Atlases}, manuscript, May 2026.
\end{thebibliography}

\end{document}
